\paragraph{Section 9 fragment statement.}
\begin{proposition}[Conjecture under test: middle facet]\label{prop:middle-facet}
Under Assumptions~\ref{ass:randomization}--\ref{ass:witness-cell}, the certificate \(C_{\mathrm{mid}}(q;\Gamma):q_b^1\le T_{\mathrm{mid}}(\Gamma)\) is a valid and irredundant facet of the projected common-law cell polytope for the endpoint coordinate \(q_b^1\).  The facet is supported by the normal \(n_{\mathrm{mid}}=e_b^1\), and its active face is attained by
\[
  q_b^0=\frac{1}{2\Gamma+1},\quad
  q_a^0=q_c^0=\frac{\Gamma}{2\Gamma+1},\quad
  q_b^1=\frac{4\Gamma-1}{3(2\Gamma+1)},\quad
  q_a^1=q_c^1=\frac{\Gamma+2}{3(2\Gamma+1)}.
\]
Moreover, for every finite \(\Gamma>1\), no pair of independent adjacent-edge optimizers attaining \(T_{\mathrm{pair}}(\Gamma)\) on both incident edges has a common middle coordinate in the symmetric witness cell, because the middle coordinate required by each edge exceeds the coherent facet endpoint by \(\Delta_{\mathrm{mid}}^{\mathrm{sym}}(\Gamma)>0\).

The middle-facet certificate is valid, irredundant in the one-dimensional endpoint projection, attained by the displayed coherent point, and separated from the incident pairwise endpoint by \(\Delta_{\mathrm{mid}}^{\mathrm{sym}}(\Gamma)>0\).
\end{proposition}

\paragraph{Section 13 refutation attempt.}
The refutation primitive is the finite vertex search for the binary-outcome, three-stratum witness cell.  Write
\[
  q_h^y=q_{z,h}(x_0,y),\qquad h\in\{a,b,c\},\quad y\in\{0,1\}.
\]
By Assumption~\ref{ass:witness-cell}, \(\Pr(Y=1\mid Z=z,S=k,X=x_0)=1/2\) and \(w_a=w_b=w_c=1/3\).  Therefore the local common-law cell polytope is the set of nonnegative arrays satisfying
\[
  q_a^y+q_b^y+q_c^y=1\quad (y=0,1),
  \qquad
  q_h^0+q_h^1=\frac23\quad (h=a,b,c),
\]
and, because the baseline weights are equal, the adjacent odds-ratio inequalities reduce to
\[
  \Gamma^{-1}q_b^y\le q_a^y\le \Gamma q_b^y,\qquad
  \Gamma^{-1}q_b^y\le q_c^y\le \Gamma q_b^y,
  \qquad y=0,1.
\]
This is a two-dimensional polytope after the four independent equality constraints are imposed.  A violating vertex would have
\[
  q_b^1>\frac{4\Gamma-1}{3(2\Gamma+1)}.
\]
Since calibration gives \(q_b^1=2/3-q_b^0\), such a vertex would require
\[
  q_b^0<\frac{1}{2\Gamma+1}.
\]
The searched vertex face cannot contain such a point: the \(y=0\) adjacent constraints imply
\[
  1=q_a^0+q_b^0+q_c^0
  \le \Gamma q_b^0+q_b^0+\Gamma q_b^0
  =(2\Gamma+1)q_b^0.
\]
Hence every feasible vertex, and therefore every feasible point, has \(q_b^0\ge (2\Gamma+1)^{-1}\).  The only candidate saturating this lower bound forces
\[
  q_a^0=q_c^0=\Gamma q_b^0=\frac{\Gamma}{2\Gamma+1},
\]
and calibration then forces exactly the displayed \(y=1\) coordinates.  The finite vertex search therefore produces no refuting witness.  The calculation is the standard extreme-point/linear-programming check for a finite partial-identification polytope, as in Manski (1990) and Balke--Pearl (1997).

\paragraph{Section 10 interpretation.}
The result says that the middle principal stratum cannot be optimized on the two incident adjacent edges independently.  The common-law calibration ties the two binary-outcome posterior rows together: reducing \(q_b^0\) raises the endpoint coordinate \(q_b^1\), but the two neighbors \(a\) and \(c\) must jointly absorb the remaining \(y=0\) mass while each stays within a \(\Gamma\)-odds-ratio band around \(b\).  The coherent K=3 endpoint is therefore determined by the shared middle-stratum bottleneck
\[
  T_{\mathrm{mid}}(\Gamma)=\frac{4\Gamma-1}{3(2\Gamma+1)}.
\]
Theorem~\ref{thm:parent-reduction} explains why this cell-level certificate is relevant to the parent sharp set: the reported ordinal-principal interval is the image of the same feasible posterior polytope under the linear target map.  The present proposition adds the extra compatibility facet that is invisible if the two incident adjacent edges are optimized as separate pairwise relaxations.

\paragraph{Section 11 proof.}
Fix the positive witness cell from Assumption~\ref{ass:witness-cell}.  All coordinates below are evaluated only in this cell, so Assumption~\ref{ass:positivity} supplies the needed support and conditioning conditions.  Assumption~\ref{ass:divergence} supplies the adjacent odds-ratio inequalities, and the equal witness weights make those inequalities ordinary ratio inequalities among the posterior coordinates.

Let
\[
  x=q_b^0,\qquad u=q_a^0,\qquad v=q_c^0.
\]
The \(y=0\) simplex constraint gives
\[
  u+x+v=1.
\]
The adjacent odds-ratio restrictions at \(y=0\) give \(u\le \Gamma x\) and \(v\le \Gamma x\).  Hence
\[
  1=u+x+v\le \Gamma x+x+\Gamma x=(2\Gamma+1)x,
\]
so
\[
  q_b^0=x\ge \frac{1}{2\Gamma+1}. \tag{1}
\]
Calibration for the middle stratum gives
\[
  \frac12 q_b^0+\frac12 q_b^1=\frac13,
\]
or equivalently
\[
  q_b^1=\frac23-q_b^0. \tag{2}
\]
Combining (1) and (2) yields the valid endpoint bound
\[
  q_b^1\le \frac23-\frac{1}{2\Gamma+1}
  =\frac{4\Gamma-1}{3(2\Gamma+1)}
  =T_{\mathrm{mid}}(\Gamma). \tag{3}
\]

Now define the displayed candidate by
\[
  q_b^0=\frac{1}{2\Gamma+1},\qquad
  q_a^0=q_c^0=\frac{\Gamma}{2\Gamma+1},
\]
and
\[
  q_b^1=\frac{4\Gamma-1}{3(2\Gamma+1)},\qquad
  q_a^1=q_c^1=\frac{\Gamma+2}{3(2\Gamma+1)}.
\]
The two simplex equations hold because
\[
  \frac{\Gamma+1+\Gamma}{2\Gamma+1}=1
\]
and
\[
  \frac{\Gamma+2+4\Gamma-1+\Gamma+2}{3(2\Gamma+1)}=1.
\]
The three calibration equations hold because, for \(h=b\),
\[
  \frac{1}{2\Gamma+1}+\frac{4\Gamma-1}{3(2\Gamma+1)}=\frac23,
\]
and, for \(h=a,c\),
\[
  \frac{\Gamma}{2\Gamma+1}+\frac{\Gamma+2}{3(2\Gamma+1)}=\frac23.
\]
The \(y=0\) adjacent inequalities hold with equality on both incident edges:
\[
  q_a^0=q_c^0=\Gamma q_b^0.
\]
For \(y=1\), the only nontrivial upper inequality is
\[
  q_b^1\le \Gamma q_a^1=\Gamma q_c^1,
\]
which is equivalent to
\[
  4\Gamma-1\le \Gamma(\Gamma+2),
\]
or \((\Gamma-1)^2\ge 0\).  The reverse lower inequalities are equivalent to
\[
  \Gamma(4\Gamma-1)\ge \Gamma+2,
\]
or \(2(2\Gamma+1)(\Gamma-1)\ge 0\).  Thus the displayed point is feasible for every finite \(\Gamma>1\) and attains equality in (3).

The support normal for (3) in the projected endpoint coordinate is \(e_b^1\), since the inequality is exactly \(\langle e_b^1,q\rangle\le T_{\mathrm{mid}}(\Gamma)\).  The projection of the common-law cell polytope onto \(q_b^1\) is the interval
\[
  \Pi_b\mathcal Q_{\mathrm{cell}}(\Gamma)
  =
  \{q_b^1:q\in\mathcal Q_{\mathrm{cell}}(\Gamma)\}.
\]
Equation (3) and the displayed attaining point show that its upper endpoint is \(T_{\mathrm{mid}}(\Gamma)\).  Therefore the exposed endpoint face is nonempty and is a facet of the one-dimensional projected polytope.  It is irredundant because deleting the inequality \(q_b^1\le T_{\mathrm{mid}}(\Gamma)\) from the endpoint description strictly enlarges the projected interval above its true upper endpoint.  In the unprojected two-dimensional cell polytope, the same support inequality exposes the unique vertex displayed above.

It remains to compare the coherent endpoint with the independent adjacent-edge optimizers.  Consider the pairwise relaxation for the incident edge \(a\sim b\): keep simplex, calibration, nonnegativity, and the \(a\)-\(b\) adjacent odds-ratio inequalities, but do not impose the \(c\)-\(b\) inequalities.  With \(x=q_b^0\) and \(u=q_a^0\), nonnegativity of \(q_c^1\) gives
\[
  q_c^1=\frac23-q_c^0
  =\frac23-(1-x-u)
  =x+u-\frac13\ge 0,
\]
so \(u\ge 1/3-x\).  The \(y=0\) edge inequality gives \(u\le \Gamma x\).  Hence
\[
  \frac13-x\le \Gamma x,
\]
and
\[
  x\ge \frac{1}{3(\Gamma+1)}.
\]
This lower bound is attained by
\[
  q_b^0=\frac{1}{3(\Gamma+1)},\qquad
  q_a^0=\frac{\Gamma}{3(\Gamma+1)},\qquad
  q_c^0=\frac23,
\]
with the calibrated \(y=1\) coordinates.  Therefore the pairwise middle endpoint is
\[
  T_{\mathrm{pair}}(\Gamma)
  =
  \frac23-\frac{1}{3(\Gamma+1)}
  =
  \frac{2\Gamma+1}{3(\Gamma+1)}.
\]
The same calculation applies symmetrically to the edge \(c\sim b\).  Thus any pair of independent incident-edge optimizers attaining their pairwise endpoint requires the common middle coordinate
\[
  q_b^1=T_{\mathrm{pair}}(\Gamma).
\]
But the coherent common-law cell polytope requires \(q_b^1\le T_{\mathrm{mid}}(\Gamma)\), and
\[
  \Delta_{\mathrm{mid}}^{\mathrm{sym}}(\Gamma)
  :=
  T_{\mathrm{pair}}(\Gamma)-T_{\mathrm{mid}}(\Gamma)
  =
  \frac{2\Gamma+1}{3(\Gamma+1)}
  -\frac{4\Gamma-1}{3(2\Gamma+1)}
  =
  \frac{\Gamma+2}{3(\Gamma+1)(2\Gamma+1)}
  >0
\]
for every finite \(\Gamma>1\).  Hence no pair of independent adjacent-edge optimizers attaining \(T_{\mathrm{pair}}(\Gamma)\) on both incident edges has a common middle coordinate in the coherent symmetric witness cell.  This proves the proposition.

\paragraph{Section 12 sanity check.}
At \(\Gamma=1\), the adjacent odds-ratio restrictions collapse to equality of posterior coordinates along each incident edge.  The coherent formula gives
\[
  q_b^0=\frac13,\qquad q_a^0=q_c^0=\frac13,\qquad
  q_b^1=\frac13,\qquad q_a^1=q_c^1=\frac13,
\]
so the common-law cell is the no-divergence posterior array.  The endpoint formula also gives
\[
  T_{\mathrm{mid}}(1)=\frac{4-1}{3(2+1)}=\frac13,
\]
which is exactly the value forced by equality of all adjacent posterior coordinates.  This is the expected no-sensitivity corner.

\paragraph{Section 13 counterexample or minimality check.}
The common-law facet is minimal in the following concrete sense.  If the \(c\)-\(b\) adjacent restriction is dropped while the \(a\)-\(b\) pairwise relaxation is retained, the pairwise optimizer above is feasible for that relaxation and has
\[
  q_b^1=T_{\mathrm{pair}}(\Gamma)>T_{\mathrm{mid}}(\Gamma).
\]
It fails the omitted coherent restriction because
\[
  \frac{q_c^0}{q_b^0}
  =
  \frac{2/3}{1/\{3(\Gamma+1)\}}
  =
  2(\Gamma+1)>\Gamma.
\]
Thus the endpoint gain is not a harmless reparameterization: it is obtained exactly by violating the missing incident-edge compatibility.  Symmetrically, dropping the \(a\)-\(b\) restriction allows the \(c\)-\(b\) pairwise optimizer and violates \(q_a^0\le \Gamma q_b^0\).  This witness explains why the shared middle-stratum constraint is necessary for the strict nonrectangularity conclusion.
